
% =====================================================================
%  Figure 4.1 -- Overview of the proposed system  (REVISED)
%  This is your own TikZ picture with the agreed corrections applied.
%  Nothing outside the picture changed: same macros (\PNBlock, \PNCaption,
%  \DNA, \Broom, ...), same styles (dashedPanel, panelTitle, card, arrow,
%  orangeArrow, redArrow, greenArrow), same colours.
%
%  WHAT CHANGED (search for "CHANGED:" to find each edit in place)
%   1. Pre-processing panel: "Cleaning + deduplication", "Normalization
%      Z=(X-mu)/sigma" and "Feature reduction: constant + correlation
%      filter" all removed -- none of the three is in the thesis.
%      Replaced by record matching, phenotype encoding and the label mask.
%   2. NEW Gene Selection band placed AFTER the split, fed by all three
%      protocols -- this is the leakage-order fix.
%   3. Data Split panel + strategy cards carry the real per-split sizes,
%      including the lineage holdout's 4,836 / 1,081 / 1,045.
%   4. Lens B is labelled strain-conditioned.
%   5. All four evaluation pillars are present and numbered; a new
%      Interpretability panel is added.
%   6. GeneAttentionNet becomes a design-lineage note; PR-AUC 0.769 gone,
%      and the misleading arrow from Lens C is removed.
%   7. Giessen: 90.4% -> 90.0%, and "same TriLens methodology" corrected
%      to the classical stacking arm only.
%   8. The three coloured split arrows now converge on one place instead
%      of running unlabelled to every block, and a legend is given.
%
%  Block contents were kept to short labels -- the detail belongs in the
%  Chapter 4 prose, not in the workflow figure. Bounding box: -44.50 ->
%  -45.20, aspect ratio 1.57 -> 1.60, so the float still fits its page.
% =====================================================================
 
\begin{tikzpicture}[x=1cm,y=1cm,scale=.82,transform shape]
\renewcommand{\Huge}{\fontsize{28.5}{34}\selectfont}
\tikzset{
  arrow/.append style={shorten >=1.2pt, rounded corners=2pt},
  orangeArrow/.append style={shorten >=1.2pt, rounded corners=2pt},
  redArrow/.append style={shorten >=1.2pt, rounded corners=2pt},
  greenArrow/.append style={shorten >=1.2pt, shorten <=0.5pt, rounded corners=2pt},
  % Small numbered "stage" badges + zone corner tags used to make the
  % top-to-bottom reading order of the redesigned pipeline explicit.
  stageBadge/.style args={#1}{circle,fill=#1,text=white,draw=black!25,
    line width=.5pt,font=\bfseries\Large,minimum size=1.05cm,inner sep=0pt},
  zoneTag/.style args={#1}{font=\bfseries\Huge,text=#1!65!black,anchor=west},
}
 
% CHANGED: bounding box nudged to fit the new gene-selection band.
\path[use as bounding box] (-1.00,-45.20) rectangle (45.70,30.30);
 
% Shift the complete pipeline upward while retaining all internal alignments.
\begin{scope}[yshift=1.5cm]
 
% =====================================================================
% Dataset panel
% =====================================================================
\begin{scope}[xshift=2.30cm,yshift=.5cm]
\node[
  dashedPanel,
  draw=navy,
  minimum width=11.6cm,
  minimum height=11.10cm
] (datasetBox) at (3.20,22.25) {};
 
\node[panelTitle,text=navy,anchor=north,font=\bfseries\Huge]
  at (3.20,27.25) {Dataset};
 
\node[font=\bfseries\Huge,align=center]
  at (3.20,25.05) {Arcadia E. coli\\AMR dataset};
 
\DNA{(0.30,19.35)}{1.8}
 
\foreach \x/\y/\c in {
  2.05/23.48/blueedge,
  2.72/23.68/orangeedge,
  3.35/23.10/navy,
  4.05/23.48/orangeedge,
  4.70/23.20/blueedge,
  5.18/23.58/orangeedge,
  1.82/22.55/orangeedge,
  2.62/22.28/blueedge,
  3.40/22.52/blueedge,
  4.12/22.60/orangeedge,
  4.78/22.35/navy,
  5.28/22.82/blueedge,
  2.12/21.72/navy,
  3.18/21.82/orangeedge,
  4.18/21.76/blueedge%
}{
  \fill[fill=\c] (\x,{\y-1.5}) circle (.10);
}
 
\Bacteria{(5.00,20.15)}{2}
 
% CHANGED: the two source deposits are flagged, because they feed two
% different views. Only "Release 2" is a name the thesis uses (Section
% 4.3.2.2); the variant-call deposit is left unnamed rather than guessed at.
\node[font=\Huge,align=center] at (3.20,17.75) {6,962 genomes};
\end{scope}
 
% =====================================================================
% Pre-processing panel
% CHANGED: this panel used to read "Cleaning + deduplication / duplicates
% removed", "Normalization Z=(X-mu)/sigma" and "Feature reduction:
% constant + correlation filter". None of those three appears anywhere in
% the thesis: Section 4.3.1 says the genomes needed no cleaning, no
% z-scoring is described, and the real filters are the two in Section
% 4.3.4 (which have now moved below the split, where they belong).
% =====================================================================
\begin{scope}[xshift=7.55cm,yshift=.5cm]
\node[
  rounded corners=2mm,
  draw=orangeedge!60,
  fill=orangefill!28,
  line width=.7pt,
  minimum width=14.5cm,
  minimum height=11.10cm
] (preBox) at (13,22.25) {};
 
\node[panelTitle,text=navy,anchor=north,font=\bfseries\Huge,align=center]
  at (13,27.25) {Preprocessing\\Framework};
 
\node[
  card=orangeedge,
  minimum width=13.5cm,
  minimum height=2.0cm
] (matchCard) at (13,23.35) {
  \bfseries\Huge{Record matching}\\[4pt]
  \Huge three files intersected
};
 
\node[
  card=orangeedge,
  minimum width=13.5cm,
  minimum height=2.0cm
] (labelCard) at (13,20.65) {
  \bfseries\Huge{Phenotype encoding}\\[4pt]
  \Huge R, I $\rightarrow$ 1 \quad S $\rightarrow$ 0
};
 
\node[
  card=orangeedge,
  minimum width=13.5cm,
  minimum height=2.0cm
] (maskCard) at (13,18.15) {
  \bfseries\Huge{Label mask}\\[4pt]
  \Huge $6{,}962\times7$
};
 
\Broom{(6.55,22.55)}{.85}
\Bars{(6.55,20.15)}{.85}
\Funnel{(6.55,17.95)}{.75}
\end{scope}
 
\draw[arrow] (datasetBox.east) -- (preBox.west);
 
% =====================================================================
% Data split panel
% =====================================================================
\begin{scope}[xshift=12.25cm,yshift=.5cm]
\node[
  rounded corners=2mm,
  draw=blueedge!60,
  fill=bluefill!48,
  line width=.7pt,
  minimum width=14.5cm,
  minimum height=11.10cm
] (splitPanelBox) at (24.8,22.25) {};

\node[panelTitle,text=navy,anchor=north,font=\bfseries\Huge,align=center]
  at (24.8,27.25) {Data Split\\Framework};

\node[
  card=blueedge,
  minimum width=12.6cm,
  minimum height=2.0cm
] (training) at (24.8,23.25) {
  \bfseries\Huge{Training (70\%)}\\[2pt]
  \Huge 4,872 genomes
};

\node[
  card=blueedge,
  minimum width=12.6cm,
  minimum height=2.0cm
] (validation) at (24.8,20.7) {
  \bfseries\Huge{Validation (15\%)}\\[2pt]
  \Huge 1,045 genomes
};

\node[
  card=blueedge,
  minimum width=12.6cm,
  minimum height=2.0cm
] (testing) at (24.8,18.2) {
  \bfseries\Huge{Testing (15\%)}\\[2pt]
  \Huge 1,045 genomes
};

\GradCap{(19,22.85)}{.85}
\Clipboard{(19,20.3)}{.85}
\Clipboard{(19,17.8)}{.75}


\end{scope}

\node[
  minimum width=12.9cm,
  minimum height=8.9cm,
  inner sep=0pt
] (splitBox) at (31.1,22.1) {};

% Data-splitting strategies
\node[
  card=blueedge,
  minimum width=6.6cm,
  minimum height=2.15cm,
  font=\bfseries\Huge
] (randomSplit) at ([xshift=-8.2cm,yshift=-3cm]splitBox.south) {
  Random Split
};

\node[
  card=blueedge,
  minimum width=7.6cm,
  minimum height=2.15cm,
  font=\bfseries\Huge
] (phylogroupSplit) at ([yshift=-3cm]splitBox.south) {
  Phylogroup-\\Stratified
};

\node[
  card=blueedge,
  minimum width=7.2cm,
  minimum height=2.15cm,
  font=\bfseries\Huge
] (lineageHoldout) at ([xshift=8.2cm,yshift=-3cm]splitBox.south) {
  Lineage\\Holdout
};

% One shared trunk from the centre of the Data Split panel's bottom edge,
% forked to the three splitting strategies. Pin the start to the panel frame
% explicitly so it cannot appear to originate from the Testing card.
\coordinate (splitStrategyStart) at
  (splitPanelBox.south);
\coordinate (splitStrategyFork) at
  ([yshift=-0.5cm]splitStrategyStart);
\draw[line] (splitStrategyStart) -- (splitStrategyFork);
\draw[line]
  (randomSplit.north |- splitStrategyFork) --
  (lineageHoldout.north |- splitStrategyFork);
\draw[arrow]
  (randomSplit.north |- splitStrategyFork) -- (randomSplit.north);
\draw[arrow]
  (phylogroupSplit.north |- splitStrategyFork) -- (phylogroupSplit.north);
\draw[arrow]
  (lineageHoldout.north |- splitStrategyFork) -- (lineageHoldout.north);

\draw[arrow] (preBox.east) -- (splitPanelBox.west |- preBox.east);
 
% =====================================================================
% CHANGED -- NEW STAGE: gene selection, AFTER the split.
%
% This is the single most important correction. The old figure put
% "Feature reduction" inside the pre-processing panel, i.e. BEFORE the
% split, which is exactly the leaked pipeline the thesis says it found and
% fixed (Section 4.3.5: "re-fitted on each split's training partition
% instead of being computed only once ... The definition of the pipeline
% used an earlier version that leaked information between splits, which
% was fixed prior to any results in this thesis").
% =====================================================================
\PNBlock{stage15}{(25.8,10.5)}{30cm}{2.70cm}{orangefill}{orangeedge}{
  \Huge\textbf{Gene Selection}\\[5pt]
  \Huge $\ge$ 5 SNPs $\rightarrow$ 13,799 genes
  \quad$\bullet$\quad
  XGBoost gain $\rightarrow$ \textbf{top 5,000 genes}
}
 
 
% Straight arrows carry the three split protocols into gene selection.
\draw[arrow]
  (randomSplit.south) -- (randomSplit.south |- stage15.north);
\draw[greenArrow]
  (phylogroupSplit.south) -- (phylogroupSplit.south |- stage15.north);
\draw[redArrow]
  (lineageHoldout.south) -- (lineageHoldout.south |- stage15.north);
 
% =====================================================================
% Genomic input blocks
% CHANGED: pushed down to make room for the gene-selection band, and now
% fed from it rather than straight from the split.
% =====================================================================
\def\genomicViewWidth{5cm}
\def\genomicViewHeight{4.15cm}
\coordinate (genomicViewRow) at ([yshift=-12.90cm]splitBox.south);
 
\PNBlock{snp}{([xshift=-8.2cm]genomicViewRow)}{\genomicViewWidth}{\genomicViewHeight}{bluefill}{blueedge}{
  \Huge\textbf{SNP}\\[3pt]
  \Huge\textbf{blocks}
}
\PNCaption{snp}{\Huge 6,962$\times$5,000$\times$64}
\DNA{([xshift=-2cm,yshift=-.62cm]snp.center)}{.60}
 
\PNBlock{gene}{(genomicViewRow)}{\genomicViewWidth}{\genomicViewHeight}{bluefill}{blueedge}{
  \Huge\textbf{Gene}\\[3pt]
  \Huge\textbf{Presence/}\\[3pt]
  \Huge\textbf{Absence}\\[3pt]
  \Huge\textbf{Matrix}
}
\PNCaption{gene}{\Huge 6,962$\times$32,441}
 
\PNBlock{summary}{([xshift=8.2cm]genomicViewRow)}{\genomicViewWidth}{\genomicViewHeight}{bluefill}{blueedge}{
  \Huge\textbf{Genomic}\\[4pt]
  \Huge\textbf{Summary}\\[4pt]
  \Huge\textbf{Stats}
}
\PNCaption{summary}{{\Huge 6,962$\times$5,000$\times$4}}
 
% Classical branch. CHANGED: its two distinct inputs are now named -- four
% learners read the gene-level features, XGB_raw reads the chi-squared
% prefiltered raw SNP matrix (Section 4.2.2, Table 4.8).
\PNBlock{ensemble}{([xshift=-8.5cm,yshift=-2.4cm]snp.west)}{11.3cm}{4.80cm}{orangefill}{orangeedge}{
  {\Huge\itshape Classical Arm}\\[3pt]
  \Huge\textbf{Classical ML Ensemble}\\[3pt]
  \Huge\textbf{5 base learners (OOF)}\\[4pt]
  \Huge XGB\_gene, LGBM, RF,\\[3pt]
  \Huge LogReg, XGB\_raw\\[3pt]
  \Huge XGB\_raw: $\chi^2$\\[3pt]
  \Huge top-50,000 raw SNPs
}
 
\node[font=\Huge,anchor=north]
  at ([yshift=-.28cm]ensemble.south) {(5-fold OOF)};
 
% CHANGED: gene selection feeds only the two views that are built from the
% selected 5,000 genes, plus the classical arm's gene-level feature matrix
% (Section 4.1, "the identical gene features feed five base learners";
% Figure 4.4 shows that matrix as 5,000 genes x 4 stats).
\draw[orangeArrow] (snp.north     |- stage15.south) -- (snp.north);
\draw[orangeArrow] (summary.north |- stage15.south) -- (summary.north);
\draw[orangeArrow] (11.9,0 |- stage15.south) -- ++(0,-0.7) -| (ensemble.north);
\node[font=\Huge,anchor=south,text=black!82]
  at (13.9,6.4) {gene-level features ($5{,}000\times4$)};
 
% CHANGED: the pan-genome view is NOT filtered -- release 2 keeps all 32,441
% loci (Section 4.3.2.2, Table 4.6), so it goes round the band, not through it.
% a line-jump keeps it unmistakable that this path ends at the pan-genome
% view only -- it merely passes over the per-gene-statistics arrow.
\draw[draw=orangeedge,line width=3.4pt,dashed,rounded corners=2pt]
  (44.30,19.20) -- (44.90,19.20) -- (44.90,7.60) -- (39.78,7.60);
\draw[draw=orangeedge,line width=3.4pt]
  (39.78,7.60) arc (0:180:0.48);
\draw[orangeArrow,dashed]
  (38.82,7.60) -- (31.10,7.60) -- (gene.north);
\node[font=\Huge,anchor=south,text=black!82]
  at (34.60,7.78) {no gene selection};
 
% CNN stack fed by the SNP blocks.
\PNBlock{cnnOne}{([xshift=-.47cm,yshift=-3.2cm]snp.south)}{.52cm}{1.75cm}{greenfill}{greenedge}{}
\PNBlock{cnnTwo}{([yshift=-3.2cm]snp.south)}{.52cm}{1.75cm}{greenfill}{greenedge}{}
\PNBlock{cnnThree}{([xshift=.47cm,yshift=-3.2cm]snp.south)}{.52cm}{1.75cm}{greenfill}{greenedge}{}
 
\node[font=\Huge,anchor=north]
  at ([yshift=-.28cm]cnnTwo.south) {16 $\rightarrow$ 32 $\rightarrow$ pool};
 
\draw[orangeArrow] (snp.south) -- (cnnTwo.north);
 
% Align the three neural modules beneath their corresponding genomic views.
\def\neuralModuleWidth{6.75cm}
\def\neuralModuleHeight{5.20cm}
\coordinate (neuralModuleRow) at ([yshift=-8.35cm]gene.south);

\PNBlock{motif}{(snp.center |- neuralModuleRow)}{\neuralModuleWidth}{\neuralModuleHeight}{greenfill}{greenedge}{
  {\huge\itshape Lens A}\\[4pt]
  \Huge\textbf{Motif Pattern}\\[3pt]
  \Huge\textbf{Encoder}\\[3pt]
  \Huge(SNP-motif CNN)
}

\PNBlock{attention}{([xshift=-.2cm]gene.center |- neuralModuleRow)}{8cm}{\neuralModuleHeight}{greenfill}{greenedge}{
  {\huge\itshape Lens B}\\[4pt]
  \Huge\textbf{Gene Feature}\\[3pt]
  \Huge\textbf{Attention}\\[3pt]
  \Huge\textbf{Module}\\[3pt]
  \Huge(Gated Attention)
}

\PNBlock{context}{(summary.center |- neuralModuleRow)}{\neuralModuleWidth}{\neuralModuleHeight}{greenfill}{greenedge}{
  {\huge\itshape Lens C}\\[4pt]
  \Huge\textbf{Cross-Antibiotic}\\[3pt]
  \Huge\textbf{Context}\\[3pt]
  \Huge\textbf{Module}\\[3pt]
  \Huge(Multi-task Head)
}

 
% Connect each genomic input to its corresponding neural module.
\draw[orangeArrow] (cnnTwo.south) -- (motif.north);
\draw[orangeArrow] (gene.south) -- ([xshift=.2cm]attention.north);
\draw[orangeArrow] (summary.south) -- (context.north);
 
% Mixture-of-Experts gating unit fed by all three neural modules.
% CHANGED: the gate weights are per sample AND per antibiotic -- without
% that, the box reads as an ordinary static three-way ensemble.
\PNBlock{gate}{([xshift=.7cm,yshift=-4.60cm]attention.south)}{9.6cm}{5.20cm}{orangefill}{orangeedge}{
  {\huge\itshape Neural Lens Fusion}\\[3pt]
  \Huge\textbf{Mixture-of-Experts}\\[3pt]
  \Huge\textbf{Gating Unit}\\[3pt]
  \Huge concat $z$ (192)\\[3pt]
  \Huge$\rightarrow$ softmax
}
 
\node[font=\Huge,anchor=north]
  at ([yshift=-.08cm]gate.south)
  {weights (A,B,C)};
 
\coordinate (gateArrowNorth) at ([xshift=-.5cm]gate.north);
\coordinate (gateArrowWest) at ([xshift=-.5cm]gate.west);
 
\draw[orangeArrow,dashed]
  (motif.south) -- ++(0,-.55) -| ([xshift=-.85cm]gateArrowNorth);
\draw[orangeArrow,dashed]
  ([xshift=.2cm]attention.south) -- (gateArrowNorth);
\draw[orangeArrow,dashed]
  (context.south) -- ++(0,-.55) -| ([xshift=.85cm]gateArrowNorth);
 
% Fuse the seven neural probabilities produced by the gated experts.
\coordinate (pipelineCenterLine) at (21,0);
\PNBlock{fusion}{(pipelineCenterLine |- gate.center)}{7.4cm}{3.45cm}{bluefill}{blueedge}{
  \Huge\textbf{Representation}\\[3pt]
  \Huge\textbf{Fusion}\\[3pt]
  \Huge neural prob (7)
}
 
\draw[orangeArrow] (gate.west) -- (fusion.east);
 
% Stacking meta-learner.
% CHANGED: the six inputs it actually receives are named.
\PNBlock{meta}{([yshift=-4.90cm]fusion.center)}{9.0cm}{4.35cm}{redfill}{rededge}{
  \Huge\textbf{Meta Learner}\\[3pt]
  \Huge(Logistic Regression)\\[3pt]
  \Huge stacking, \emph{per antibiotic}\\[3pt]
  \Huge(neural + classical)
}
 
\draw[orangeArrow] (fusion.south) -- (meta.north);
 
% CHANGED: probability and hard call are separated -- the sigmoid produces
% probabilities (which is what PR-AUC is computed on) and the threshold
% turns them into R/S.
\PNBlock{threshold}{([yshift=-2.50cm]meta.south)}{6.6cm}{3.00cm}{grayfill}{black!48}{
  \Huge\textbf{Sigmoid}\\[3pt]
  \Huge 7 probabilities\\[3pt]
  \Huge threshold 0.5
}
 
\draw[orangeArrow] (meta.south) -- (threshold.north);
 
% Final multi-label output for the seven antibiotics.
\coordinate (outputsBoxCenter) at ($(threshold.south)+(0,-5cm)$);
\node[
  rounded corners=2mm,
  draw=rededge!70,
  fill=redfill!18,
  line width=.8pt,
  minimum width=17cm,
  minimum height=8cm
] (outputsBox) at (outputsBoxCenter) {};
 
\node[
  font=\bfseries\Huge,
  text=rededge,
  anchor=north
] at ([yshift=-.36cm]outputsBox.north) {
  Output: 7 antibiotics (multi-label)
};
 
\foreach \nm/\xx/\yy in {
  AMP/-5.70/1.30,
  CIP/-1.90/1.30,
  AMC/1.90/1.30,
  SXT/5.70/1.30,
  GEN/-3.80/-1.80,
  CTX/0/-1.80,
  CAZ/3.80/-1.80
}{
  \PNBlock{out\nm}{([xshift=\xx cm,yshift=\yy cm]outputsBox.center)}{3.0cm}{1.25cm}{redfill}{rededge}{
    \Huge\textbf{\nm}
  }
 
  \node[font=\Huge,anchor=north]
    at ([yshift=-.18cm]out\nm.south) {R / S};
}
 
\draw[redArrow] (threshold.south) -- (outputsBox.north);
 
% Feed the classical ensemble into the representation-fusion block.
\draw[orangeArrow]
  (ensemble.south) |- (fusion.west);
 
% =====================================================================
% Pillar 1 -- Locality Test.
% CHANGED: moved down (the classical arm now occupies its old row), given
% its pillar number, renamed so it can no longer be confused with Lens A's
% SNP-motif CNN, and its claim softened to what the experiment supports.
% =====================================================================
\node[
  dashedPanel,
  draw=greenedge,
  minimum width=9.5cm,
  minimum height=15.90cm
] (localityBox) at (3.95,-12.60) {};
 
\node[panelTitle,text=greenedge,anchor=north,font=\bfseries\Huge,align=center]
  at ([yshift=-.35cm]localityBox.north) {Pillar 1\\Locality Test};
 
\PNBlock{loccnn}{([yshift=3.60cm]localityBox.center)}{5.80cm}{2.30cm}{greenfill}{greenedge}{
  \bfseries\Huge{Locality probe}\\[3pt]
  \bfseries\Huge{(plain 1-D CNN)}
}
 
\node[
  anchor=north,
  align=center,
  font=\Huge,
  text width=7.4cm
] at ([yshift=-.55cm]loccnn.south) {
  Compares real SNP windows against shuffled positions to check whether the
  CNN uses position-local structure.
};
 
\draw[arrow]
  (datasetBox.south -| localityBox.north) -- (localityBox.north);
 
% =====================================================================
% Pillar 3 -- Interpretability.
% CHANGED: NEW panel. The caption promises four pillars and the old figure
% showed two; this is the one the thesis leans on most heavily
% (Section 5.6.1) and it was entirely absent.
% =====================================================================
\node[
  dashedPanel,
  draw=greenedge,
  minimum width=8.50cm,
  minimum height=8.20cm
] (interpBox) at ([xshift=12.4cm,yshift=-5.6cm]meta.east) {};
 
\node[panelTitle,text=greenedge,anchor=north,font=\bfseries\Huge,align=center]
  at ([yshift=-.35cm]interpBox.north) {Pillar 3\\Interpretability};
 
\node[
  anchor=north,
  align=center,
  font=\Huge,
  text width=7.4cm
] at ([yshift=-3.4cm]interpBox.north) {
  gate weights\\[6pt]
  gene attention\\[6pt]
  SHAP\\[6pt]
};
 
% Straight-segment S-shaped link from the meta learner to Pillar 3.
\coordinate (interpLinkBend) at ($(meta.east)!0.5!(interpBox.west)$);
\draw[greenArrow,dashed]
  (meta.east) -- (interpLinkBend |- meta.east)
  -- (interpLinkBend |- interpBox.west) -- (interpBox.west);
 
% =====================================================================
% Pillar 4 -- independent collection.
% CHANGED: "same TriLens methodology applied to an independent dataset" is
% not what happened -- Table 5.12's own caption says neither column is the
% full TriLens model, and Section 5.5.4 says the three lenses cannot be
% run on Giessen at all. Retention corrected 90.4% -> 90.0% (Table 5.12).
% =====================================================================
\node[
  dashedPanel,
  draw=tealedge,
  minimum width=25.00cm,
  minimum height=9.40cm
] (externalBox) at ([yshift=-6.0cm]outputsBox.south) {};
 
\draw[redArrow] (outputsBox.south) -- (externalBox.north);
 
\node[
  font=\bfseries\Huge,
  text=tealedge,
  anchor=north,
  align=center,
  text width=16.25cm
] at ([yshift=-.45cm]externalBox.north) {
  Pillar 4\\
  Cross-collection Robustness
};
 
\PNBlock{giessen}{([xshift=-3.3cm,yshift=-.15cm]externalBox.center)}{5.60cm}{3.60cm}{bluefill}{blueedge}{
  \Huge\textbf{Giessen}\\[3pt]
  \Huge\textbf{E. coli}\\[4pt]
  \Huge 809 strains,\\[2pt]
  \Huge 4 drugs
}
 
\PNBlock{retention}{([xshift=3.3cm,yshift=-.15cm]externalBox.center)}{5.60cm}{3.60cm}{greenfill}{greenedge}{
  \Huge\textbf{90.0\%}\\[3pt]
  \Huge\textbf{PR-AUC}\\[3pt]
  \Huge\textbf{retention}
}
 
\draw[arrow] (giessen.east) -- (retention.west);
 
\node[
  font=\Huge,
  align=center
] at ([yshift=1.30cm]externalBox.south) {
  the \textbf{classical stacking arm only}, retrained from scratch\\[3pt]
   a method-level check, not model transfer
};
 
% =====================================================================
% Stage-flow guide: tinted zone backgrounds + numbered badges.
% =====================================================================
\begin{scope}[on background layer]
  \coordinate (neuralLensBottom) at ([yshift=-.7cm]gate.south);
  \node[rounded corners=4mm,draw=greenedge!55,fill=greenfill!22,
    line width=1pt,inner sep=16pt,
    fit=(cnnOne)(cnnTwo)(cnnThree)(motif)(attention)(context)(gate)(neuralLensBottom)]
    (neuralLensZone) {};
  \coordinate (classicalLensBottom) at ([yshift=-1cm]ensemble.south);
  \node[rounded corners=4mm,draw=orangeedge!55,fill=orangefill!22,
    line width=1pt,inner sep=14pt,
    fit=(ensemble)(classicalLensBottom)] (classicalLensZone) {};
  \node[rounded corners=4mm,draw=blueedge!55,fill=bluefill!20,
    line width=1pt,inner sep=16pt,
    fit=(fusion)(meta)(threshold)] (decisionZone) {};
\end{scope}
 
\node[stageBadge=navy] at ([xshift=5mm,yshift=6mm]datasetBox.north west) {1};
\node[stageBadge=orangeedge] at ([xshift=5mm,yshift=6mm]preBox.north west) {2};
\node[stageBadge=blueedge] at ([xshift=5mm,yshift=6mm]splitPanelBox.north west) {3};
 
% CHANGED: Stage 1.5 gets its own badge, because it is now its own stage.
\node[stageBadge=orangeedge] at ([xshift=-8mm]stage15.west) {4};
\node[stageBadge=orangeedge] (badge5) at ($(ensemble.east)!0.5!(snp.west)$) {5};
\node[stageBadge=blueedge] at ([xshift=-1.7cm,yshift=9mm]gene.north) {6};
 
\node[stageBadge=blueedge] (badge7) at ([xshift=5mm,yshift=8mm]decisionZone.north west) {7};
\node[zoneTag=blueedge] at ([xshift=3mm]badge7.east) {Fusion + Stacking};
 
% CHANGED: Pillar 2 is the architecture itself, so it is tagged on the
% model zone rather than being left unlabelled.
\node[font=\Huge,text=greenedge!70!black,anchor=north east,align=right]
  at ([xshift=-2mm]neuralLensZone.south east) {
  Pillar 2 the architecture itself\\[2pt]
};
 
\node[stageBadge=rededge] at ([xshift=-6mm,yshift=6mm]outputsBox.north west) {8};
\node[stageBadge=tealedge] at ([xshift=-1cm,yshift=-1cm]externalBox.north west) {9};
 
\end{scope}
 
\end{tikzpicture}